\chapter{Conclusions and Future Work}
\label{chap:ch5}

\tolerance=9999
\emergencystretch=3em

\indent \par This thesis began with an awkward observation: on a cloud-hosted streaming service the bytes of a song are paid for twice, and the second payment, the egress out of the data center, grows with the number of people listening rather than with the revenue they bring. The more a track is loved, the more it costs to serve. Sitting unused inside that problem was a delivery network made of the listeners themselves, since anyone who has just heard a piece of a song holds a copy that the next listener is about to need. The question the rest of the thesis set out to answer was whether that latent network could be put to work without the listener ever noticing, which meant satisfying three constraints at once: playback had to start instantly, the same code had to run in a browser and on a phone, and a piece of audio handed over by an untrusted peer had to be exactly as safe as one fetched from the server.

\par On the evidence gathered in Chapter \ref{chap:ch4}, the answer is yes, with the caveats that a prototype always carries. The chapters between the question and this one built the argument, the system and the measurement in turn, and they line up against the three constraints and the cost goal as follows.

\indent \par The cost goal was the reason for the whole exercise, and it is the one the measurements speak to most directly. Chapter \ref{chap:ch2} derived an offload ceiling of $\rho = 1 - P/S$ and argued that a peer-assisted swarm moves the origin away from serving $O(N)$ complete streams of the same song. Chapter \ref{chap:ch4} measured a warm two-to-five-client swarm serving 88.4\% of its bytes from peers and only 11.6\% from the origin, and, once a signaling concurrency defect was found and fixed, that server share stayed flat as the swarm grew from one listener to five. A flat share is the concrete form of the claim in this prototype: the origin's involvement in any one listener's stream does not grow with the size of the audience. The figure lands in the same regime as the 8.8\% that Kreitz and Niemel\"a measured for Spotify's production system \cite{spotify_p2p}, which is the closest external point of reference that exists, and the gap between the two is accounted for by the cold cache of the prototype rather than by any difference in the mechanism.

\par The instant-playback constraint was met by the prefix. A dynamic prefix of about 5\% of each song, served from the origin, hides the WebRTC connection setup, which the measurements put at a median of 243 milliseconds, comfortably inside the eleven seconds of audio that the prefix buys. Click-to-sound therefore stays at the level of an ordinary server stream even though the bulk of the song arrives over a peer connection that did not exist when the user pressed play. The cross-platform constraint was met by building the client once, in Flutter, and running it unchanged on Linux, macOS, Windows, Android, iOS and, through a Service Worker, the web browser. The integrity constraint was met by content-addressing every 64 KiB chunk with its SHA-256 hash, so that a chunk is verified identically whether it came from a peer or the origin, and a failed verification simply falls back to the server. That fallback, rather than a TURN relay, is also what the system uses when a direct peer connection cannot be established, which keeps the cost-saving property of the peer model intact: a relayed byte would be billed exactly like a server byte.

\par It is worth being equally clear about what was not shown. The measurements were taken on a single machine over the loopback interface, with one synthetic track and listeners that always started cold, so the offload and the setup latency are best-case figures and the absolute savings are too small to quote in money. The cache share that dominates a mature deployment was zero by construction. The evaluation is therefore a demonstration that the mechanism behaves as the model predicts under controlled conditions, not a measurement of a production service. What it does establish, and what the introduction asked for, is that the latent delivery network can be put to work, that the origin can be reduced from serving the whole song to serving a small bootstrap and fallback role, and that the price of doing so, paid in a small server prefix and a few hundred milliseconds of background connection setup, need not be visible to the listener at all.

\indent \par The most valuable next step is also the most obvious one: repeat the evaluation outside the laboratory. Running the swarm across several machines on a real wide-area network, with non-trivial round-trip times and genuine NAT traversal, and over a realistic catalogue of varied music rather than a single tone, would turn the share-based results of Chapter \ref{chap:ch4} into a saving that can be quoted in currency and would expose the system to the cross-platform connection failures that the loopback interface never provokes. Warm-cache and larger-swarm runs belong in the same effort, since the cache share that this thesis measured as zero is precisely the contribution that grows to dominate a mature peer-assisted service \cite{spotify_p2p}, and the flash-crowd behaviour at large $N$, which the signaling fix addressed only at the small scale tested here, still needs to be characterised.

\par A second strand concerns resilience and reach. The cross-platform connectivity limitation of Section \ref{sec:limitations}, where browsers hide their local addresses behind mDNS hostnames that native peers cannot always resolve, is the clearest barrier to mixed browser-and-mobile swarms offloading as well as desktop pairs, and is worth attacking directly. The signaling WebSocket should reconnect on its own with a backoff loop rather than waiting for the next user action, so that a device is never silently absent from the swarm. On mobile, where seeding spends battery and metered data, the simple network-mode and monthly-cap controls implemented here could grow into a policy informed by charging state and connection type, and possibly into an incentive scheme, building on the prior characterisations of mobile peer-assisted streaming and of device-to-device communication \cite{spotify_mobile_p2p,d2d_survey}.

\par A third strand is about turning design arguments into further measurements and about the product itself. The prefix size and the peer timeouts were chosen on the strength of the cost model and fixed in code. Making them dynamic or exposing them as parameters would let the trade-offs that Chapter \ref{chap:ch2} reasoned about be swept and measured rather than argued. Beyond performance, a peer learns something about what its neighbours are listening to when it serves them chunks, so a privacy-preserving request scheme is a natural and important extension, and a relay tier could be offered, as a deliberate paid option rather than a default, to the small fraction of users for whom direct connectivity never succeeds \cite{relays_turn}.

\par A fourth, more ambitious strand is interoperability with other music streaming services and local media players. In the current prototype, cached audio is stored as separate verified chunks rather than as a normal playable audio file, so a song reconstructed from the cache cannot be opened directly by another application. Future work could explore a reconstruction layer that joins the cached chunks back into a conventional media file when enough of them are available, while still preserving the chunked layout needed for peer delivery. This would make the same stored data useful both inside MP33r and outside it, but it would have to be designed carefully so that file reconstruction, local media scanning and direct-file playback do not slow down the streaming path.

\par None of these changes the shape of the result. They widen the conditions under which it has been shown to hold. The core finding stands: the listeners of a song are, collectively, most of the infrastructure needed to deliver it, and a system can draw on that infrastructure while keeping playback instant, the catalogue safe and the code single-sourced across platforms. The opening worry was that the more a song is loved, the more it costs to serve. The system built here turns the same fact into an asset, because the more a song is loved, the more listeners there are holding a copy of it and the more hands are there to pass it along.
